% !TEX root = ../../main.tex
%% Défi du chapitre 9 — Carrés magiques et carré de Lo Shu
%% Résultats vérifiés par verifications/defis_06_10.py.

\rubrique{Défi}

\begin{defi}{Carrés magiques et carré de Lo Shu}{1 h 30}
Une matrice $M=(m_{ij})\in\mathcal{M}_3(\R)$ est \emph{magique de somme $s$}
si ses trois lignes, ses trois colonnes et ses deux diagonales ont toutes
pour somme $s$ ; elle est \emph{semi-magique} si l'on ne demande la
condition qu'aux lignes et aux colonnes. On note $J$ la matrice dont tous
les coefficients valent $1$.

\partieExo{A}{Résoudre le système des carrés magiques}
\begin{questions}
  \item Montrer qu'une combinaison linéaire de matrices magiques est
    magique, et que $J$ est magique de somme $3$.
  \item Soit $M$ magique de somme $s$. En additionnant la ligne et la
    colonne du milieu et les deux diagonales, montrer que $m_{22}=s/3$.
  \item Pour $(c,a,b)\in\R^{3}$, on pose
    \[
    L(c,a,b)=\begin{pmatrix}c-a&c+a+b&c-b\\ c+a-b&c&c-a+b\\ c+b&c-a-b&c+a\end{pmatrix}.
    \]
    Vérifier que $L(c,a,b)$ est magique de somme $3c$.
  \item Réciproquement, soit $M$ magique de somme $s$ ; on pose
    $c=m_{22}$, $a=c-m_{11}$ et $b=c-m_{13}$. En utilisant successivement
    les équations du système (diagonales, première ligne, première
    colonne\dots), montrer que $M=L(c,a,b)$. Combien de paramètres
    décrivent les matrices magiques de $\mathcal{M}_3(\R)$ ?
\end{questions}

\partieExo{B}{Le carré de Lo Shu}
\begin{questions}[resume]
  \item On cherche les carrés magiques dont les coefficients sont
    exactement les entiers $1,2,\dots,9$. Montrer que $s=15$ et $c=5$, puis
    que $\ens{|a|,|b|}=\ens{1,3}$. En déduire qu'il y a exactement $8$ tels
    carrés, et que $L(5,1,3)$ est le carré de Lo Shu de la figure.
  \item Vérifier que $L(c,a,b)^{\mathsf T}=L(c,a,-b)$. Expliquer pourquoi les
    huit carrés trouvés sont les images du Lo Shu par les huit symétries du
    carré (défi du chapitre 8).
\end{questions}

\medskip
\begin{center}
\begin{tikzpicture}[line cap=round]
  \foreach \v [count=\k from 0] in {4,9,2,3,5,7,8,1,6}{%
    \pgfmathtruncatemacro{\ii}{div(\k,3)}\pgfmathtruncatemacro{\jj}{mod(\k,3)}
    \ifnum\v=5\fill[otGreen!20] (\jj*1.1,-\ii*1.1) rectangle ++(1.1,-1.1);\fi
    \node[font=\large,text=otInk] at (\jj*1.1+0.55,-\ii*1.1-0.55) {$\v$};}
  \draw[otNavy,line width=0.9pt] (0,0) rectangle (3.3,-3.3);
  \draw[otNavy,line width=0.5pt] (1.1,0) -- (1.1,-3.3) (2.2,0) -- (2.2,-3.3)
    (0,-1.1) -- (3.3,-1.1) (0,-2.2) -- (3.3,-2.2);
  \begin{scope}[xshift=5.2cm]
    \foreach \v [count=\k from 0] in {4,9,2,3,5,7,8,1,6}{%
      \pgfmathtruncatemacro{\ii}{div(\k,3)}\pgfmathtruncatemacro{\jj}{mod(\k,3)}
      \pgfmathsetmacro{\cx}{\jj*1.1+0.55}\pgfmathsetmacro{\cy}{-\ii*1.1-0.55}
      \pgfmathtruncatemacro{\impair}{mod(\v,2)}
      \foreach \d in {1,...,\v}{%
        \pgfmathsetmacro{\rr}{\v>1 ? 0.34 : 0}
        \pgfmathsetmacro{\aa}{90+360*(\d-1)/\v}
        \ifnum\impair=1
          \filldraw[fill=white,draw=otInk,line width=0.5pt] ($(\cx,\cy)+(\aa:\rr)$) circle (2.5pt);
        \else
          \fill[otInk] ($(\cx,\cy)+(\aa:\rr)$) circle (2.5pt);
        \fi}}
    \draw[otNavy!40,line width=0.4pt] (0,0) rectangle (3.3,-3.3);
  \end{scope}
\end{tikzpicture}\\[0.5mm]
\legende{Le carré de Lo Shu, connu en Chine depuis l'Antiquité : chaque
ligne, colonne et diagonale a pour somme $15$. À droite, sa représentation
traditionnelle (impairs en blanc, pairs en noir).}
\end{center}

\partieExo{C}{Calcul matriciel}
\begin{questions}[resume]
  \item Montrer que $M$ est semi-magique de somme $s$ si et seulement si
    $MJ=JM=sJ$. En déduire qu'un produit de matrices semi-magiques de
    sommes $s$ et $t$ est semi-magique de somme $st$, et que l'inverse
    d'une matrice semi-magique inversible de somme $s$ est semi-magique de
    somme $1/s$ (on justifiera que $s\neq0$).
  \item Soit $L=L(5,1,3)$ le Lo Shu. Calculer $L^{2}$ : est-elle
    semi-magique ? magique ?
  \item Vérifier que $L$ est inversible d'inverse
    \[
    L^{-1}=\frac{1}{360}\begin{pmatrix}23&-52&53\\ 38&8&-22\\ -37&68&-7\end{pmatrix},
    \]
    et que $L^{-1}$ est magique : de quelle somme ?
\end{questions}
\end{defi}
